% ===================================================================
%  तक्ता क्रमांक ४ — प्रौढपण आणि म्हातारपण.
%
%  Traduction, faite en 2026, en marathi d'aujourd'hui, sur l'ido de
%  texto/io. Regles de la colonne : voir 10-takta-01.tex. Deux
%  scenes, deux numerotations : les cles portent c1 et c2.
%
%  LA NOCE OBLIGE LE MARATHI A ETRE PRECIS LA OU LE FRANCAIS RESTE
%  VAGUE, ET C'EST LE MEME EMBARRAS QUE LE VIETNAMIEN AU MEME TABLEAU.
%  « Beau-pere », « belle-mere », « beau-frere » ne disent pas de quel
%  cote ; le marathi a un mot par relation, et il faut lire la planche
%  pour savoir lequel :
%
%    सासरा (20)   le pere du conjoint
%    सासू (21)    la mere du conjoint
%    सून          l'epouse du fils
%    जावई         l'epoux de la fille
%    नणंद         la sœur du mari
%    मेहुणा (30)  le frere de l'epouse
%    आजी (23, 29) la grand-mere
%    आजोबा (33)   le grand-pere
%    पणजोबा       l'arriere-grand-pere
%    नातू (2) / नात (3)  le petit-fils, la petite-fille
%
%  ET DEUX ROLES DE LA CEREMONIE ONT LEUR NOM PROPRE : करवला (25),
%  celui qui accompagne le marie, et करवली (26), celle qui accompagne
%  la mariee. Le livret dit « compagnon d'honneur » et « demoiselle
%  d'honneur » ; le marathi n'a pas eu a composer.
%
%  LE RESTE DE LA MAISON :
%
%    वाडा (1)      la demeure — le grand logis maharashtrien
%    टेकडी (3)     la colline   हिन्दी : पहाड़ी
%    व्हरांडा (8)  la veranda
%    शेवंती (9)    les chrysanthemes
%    थाळी (15)     le plat
%    सुरी (13)     le couteau
%    आरामखुर्ची (34) le fauteuil
%    शेकोटी (29)   la cheminee
%    वाटी (15)     le bol
%    काढा          la tisane
%    मांजर (1)     le chat     हिन्दी : बिल्ली
%    शेपूट         la queue
%    खिसा (9)      la poche
%    तसबीर (44)    le portrait encadre
%    दिनदर्शिका (43) le calendrier
%
%  LE MARATHI POSTPOSE, ET LE TABLEAU L'A RAPPELE QUATRE FOIS. Voir
%  les blocs signales dans le journal : le remede reste celui du
%  tableau 3.
% ===================================================================

%%K t04-apar-1 apar
\VUsaut{4.0mm}
\VUtitre{15.0pt}{60}{तक्ता क्रमांक ४}
\VUsaut{1.6mm}
\VUtitre{11.4pt}{0}{प्रौढपण आणि म्हातारपण.}
\VUsaut{3.2mm}

%%K t04-tit-1 sub
\VUcentre{11.4pt}{0}{\textit{पहिला प्रवेश.}}
\VUsaut{1.6mm}
\VUtitre{11.4pt}{0}{लग्नाची मेजवानी.}
\VUsaut{2.6mm}

%%K t04-tit-2 sub
\VUpk{10.2pt}{40}{शरद ऋतू.}
\VUsaut{3.0mm}

%%K t04-c1-01-1 p
१. --- तरुण मोठे झाले आहेत. त्यांतला एक, जॉन, लग्न करतो आहे. आपण
\VUgras{लग्नाच्या मेजवानीला} हजर आहोत, जी एकाच टेबलाभोवती नवरा-नवरीची दोन्ही
कुटुंबे जमवते.

%%K t04-c1-02-1 p
२. --- आपण \VUgras{शरद ऋतूत} आहोत, \VUgras{सप्टेंबर महिन्यात}.
\VUgras{ऑक्टोबरचा} आणि \VUgras{नोव्हेंबरचा} थंड वारा अजून ग्रामीण भागाची हिरवी
पालवी उतरवून घेऊन गेलेला नाही. संध्याकाळची हवा कोमट आणि सुगंधी आहे; म्हणूनच
जेवण बागेकडच्या \VUgras{व्हरांड्यात} \textsuperscript{(8)} होते आहे.

%%K t04-c1-03-1 p
३. --- लांबवर \VUgras{टेकडीवर} \textsuperscript{(3)} जॉनच्या आईवडिलांचे घर आहे ---
हिरवाईत लपलेला एक सुरेख \VUgras{वाडा} \textsuperscript{(1)}. उत्तम मद्य देणारे
सुंदर द्राक्षमळे टेकडीचा उतार झाकतात. द्राक्षशेतकरी त्यांची लागवड करतो आणि
काळजी घेतो.

%%K t04-c1-04-1 p
४. --- द्राक्षमळ्यांवरून \VUgras{तितरांचा} \textsuperscript{(2)} एक थवा उडून जातो
आहे --- \VUgras{राखणदार} \textsuperscript{(5)} जवळ आल्यामुळे तो घाबरला आहे. तो
राखणदार, काखेत \VUgras{बंदूक} घेऊन आणि खांद्याला \VUgras{शिकारीची पिशवी} अडकवून,
\VUgras{झुडपे} \textsuperscript{(4)} धुंडाळतो आहे --- आपल्या
\VUgras{कुत्र्यासह} \textsuperscript{(6)}.

%%K t04-c1-05-1 p
५. --- \VUgras{व्हरांड्याच्या} बाहेर एक \VUgras{द्राक्षवेल} चढते आहे, जिच्यावरून
दाट \VUgras{घड} \textsuperscript{(7)} लोंबत आहेत.

%%K t04-c1-06-1 p
६. --- \VUgras{टेबल} \textsuperscript{(16)} सजवणारी सुंदर शरद-फुले म्हणजे
\VUgras{शेवंती} \textsuperscript{(9)}; नवरीच्या \VUgras{वडिलांनी}
\textsuperscript{(10)} त्यांतले एक आपल्या फ्रॉक कोटाच्या काजात खोवले आहे.

%%K t04-c1-07-1 p
७. --- जेवण संपत आले आहे. \VUgras{नोकर} \textsuperscript{(11)} गोड पदार्थांची ताटे
आणू लागला आहे आणि आता तो लागत नसलेली भांडी उचलणार आहे --- \VUgras{चमचे}
\textsuperscript{(14)}, \VUgras{काटे} \textsuperscript{(12)}, \VUgras{सुऱ्या}
\textsuperscript{(13)}, \VUgras{मोठ्या थाळ्या} \textsuperscript{(15)}.

%%K t04-tit-3 sub
\VUpk{10.2pt}{40}{नातेवाईक.}
\VUsaut{3.0mm}

%%K t04-c1-08-1 p
८. --- सगळे आनंदी दिसतात, आणि नवरदेवाची \VUgras{आई} \textsuperscript{(17)}
ज्या स्मिताने आपल्या मुलांकडे पाहते, ते तिच्या सुखाची साक्ष देते.

%%K t04-c1-09-1 p
९. --- हे लग्न प्रदेशातली दोन श्रीमंत कुटुंबे जोडते. जॉन, \VUgras{नवरदेव}
\textsuperscript{(18)}, मोठ्या जमीनदाराचा \VUgras{मुलगा} आहे आणि तो कुशल
कारखानदाराचा \VUgras{जावई} होतो आहे.

%%K t04-c1-10-1 p
१०. --- \VUgras{नवरा-नवरी} तरुण आहेत, तरीही त्यांचा \VUgras{वाङ्निश्चय} फार
पूर्वी झाला होता. \VUgras{नवरी} \textsuperscript{(19)}, मार्था, संत्र्याच्या
फुलांनी सजवलेल्या आपल्या \VUgras{पांढऱ्या झग्यात} मोहक दिसते आहे.

%%K t04-c1-11-1 p
११. --- तिचे \VUgras{सासरे} \textsuperscript{(20)} अशी सुशील \VUgras{सून}
मिळाल्याने फार सुखी आहेत, आणि जॉनची \VUgras{सासू} \textsuperscript{(21)} आपली
\VUgras{मुलगी} इतकी सुंदर दिसते आहे हे पाहून हसते आहे.

%%K t04-c1-12-1 p
१२. --- टेबलाच्या टोकाला उरलेले \VUgras{पाहुणे} \textsuperscript{(22)} बसले
आहेत --- \VUgras{नातेवाईक, मित्र, चुलतभाऊ, चुलतबहिणी}. काखेत रुमाल घेऊन
\VUgras{वाढप्यांचा प्रमुख} \textsuperscript{(23)} त्यांना तत्परतेने वाढतो आहे.

%%K t04-tit-4 sub
\VUpk{10.2pt}{40}{मित्रमंडळी.}
\VUsaut{3.0mm}

%%K t04-c1-13-1 p
१३. --- टेबलाच्या उजव्या बाजूला, अहो, एक \VUgras{तरुणी} \textsuperscript{(24)},
नवरीची \VUgras{मैत्रीण}, आणि मग

%%K t04-c1-13-1 p suite
तिचा \VUgras{भाऊ}, जो तिचा \VUgras{करवला} \textsuperscript{(25)} आहे. तो आपल्या
\VUgras{करवलीशी} \textsuperscript{(26)} गप्पा मारतो आहे --- ती नवरदेवाची बहीण आहे
आणि या लग्नाने मार्थाची \VUgras{नणंद} होते आहे. ती मोहक तरुणी आहे. तिच्याजवळ
सुंदर \VUgras{दागिने}, \VUgras{मोत्यांचा हार} आणि सोन्याचे \VUgras{कडे} आहेत.

%%K t04-c1-14-1 p
१४. --- त्यांच्याजवळ उभे राहून ग्लास उंचावणारे गृहस्थ कुटुंबाचे \VUgras{मित्र}
\textsuperscript{(27)} आहेत. ते नवरदेवाच्या \VUgras{साक्षीदारांपैकी} एक आहेत. ते
\VUgras{शुभेच्छा देतात}, नवरा-नवरीच्या \VUgras{तब्येतीसाठी} पितात आणि त्यांना
दीर्घ सुखाची इच्छा करतात. त्यांच्या अंगावर समारंभाचा पोशाख आहे --- फ्रॉक कोट
आणि \VUgras{चकचकीत बूट} \textsuperscript{(28)}. ते \VUgras{आजींना}
\textsuperscript{(29)} सोबत करत आहेत, ज्या \VUgras{विधवा} आहेत.

%%K t04-c1-15-1 p
१५. --- बारीक काळ्या मिशीचा, \VUgras{फ्रॉक कोटातला} \textsuperscript{(31)} तरुण
जॉनचा \VUgras{मेहुणा} \textsuperscript{(30)} आहे. तो नामवंत अभियंता आहे. तो एका फार
सुरेख \VUgras{तरुण बाईचा} \textsuperscript{(33)} शेजारी आहे --- ती मार्थाची
\VUgras{मोठी बहीण} आहे आणि त्या गोड मुलीची आई आहे, जी आपल्या चुलतभावाचा हात धरून
फूल द्यायला आणि \VUgras{शुभ संध्याकाळ} म्हणायला येते आहे. त्या बाईंच्या कानात
फार सुंदर \VUgras{कर्णफुले} आहेत आणि डोक्यावर सुरेख \VUgras{शिरोभूषण} आहे.

%%K t04-c1-16-1 p
१६. --- आपल्या \VUgras{ब्लाउजमुळे} \textsuperscript{(36)} आणि आपल्या
\VUgras{फुगीर विजारीमुळे} \textsuperscript{(37)} इतका विचित्र दिसणारा तो लहान
चुलतभाऊ फार दयनीय आहे. तो \VUgras{अनाथ} \textsuperscript{(35)} आहे. अगदी लहानपणीच
त्याने आपले वडील आणि आपली आई गमावली. त्याचे काका त्याचे \VUgras{पालक}
\textsuperscript{(38)} आहेत, आणि त्या बिचाऱ्या मुलाला वाढवण्यासाठी ते अविवाहित
राहिले. ते भला माणूस आहेत. ते जरा टकले आहेत आणि त्यांना फार कमी दिसते; ते
\VUgras{नाकावरचा चष्मा} वापरतात.

%%K t04-tit-5 sub
\VUsaut{2.4mm}
\VUpk{10.2pt}{40}{गोड पदार्थ.}
\VUsaut{3.0mm}

%%K t04-c1-17-1 p
१७. --- मेजवानीच्या मंडळींच्या मागे, \VUgras{मेजपोशाने} झाकलेल्या एका टेबलावर,
\VUgras{गोड पदार्थ} \textsuperscript{(39)} तयार आहेत.

%%K t04-c1-17-1 p suite
मोठ्या भांड्यात, अहो, फळे --- \VUgras{पीच} \textsuperscript{(40)},
\VUgras{पेअर} \textsuperscript{(41)}, \VUgras{सफरचंद} \textsuperscript{(42)},
\VUgras{जर्दाळू} \textsuperscript{(43)} आणि \VUgras{द्राक्षे}
\textsuperscript{(44)}. जवळच \VUgras{अननस} \textsuperscript{(45)}, सुंदर
\VUgras{केक} \textsuperscript{(46)}, \VUgras{मद्याच्या बाटल्या}
\textsuperscript{(48)} आणि \VUgras{शॅम्पेन} \textsuperscript{(49)}, आणि स्वच्छ
\VUgras{ताटांचे} \textsuperscript{(47)} थर.

%%K t04-c1-18-1 p
१८. --- \VUgras{वाढपी} \textsuperscript{(50)} जुन्या मद्याची \VUgras{बाटली}
\textsuperscript{(52)} \VUgras{बुचकाढणीने} \textsuperscript{(51)} उघडतो आहे.
\VUgras{बूच} \textsuperscript{(53)} फार अवघड निघेल असे दिसते. त्या वाढप्याच्या
काखेत \VUgras{रुमाल} \textsuperscript{(54)} आहे.

%%K t04-c1-19-1 p
१९. --- जेवणानंतर गृहस्थ धूम्रपानाच्या खोलीत जातील, आणि बाया नाचासाठी तयार
व्हायला जातील.

%%K t04-tit-6 sub
\VUsaut{5.0mm}
\VUcentre{11.4pt}{0}{\textit{दुसरा प्रवेश.}}
\VUsaut{1.6mm}
\VUpk{11.4pt}{40}{आजोबांचा वाढदिवस.}
\VUsaut{2.6mm}

%%K t04-tit-7 sub
\VUpk{10.2pt}{40}{भेट.}
\VUsaut{3.0mm}

%%K t04-c2-01-1 p
१. --- दहा वर्षे गेली; जॉनचे आईवडील म्हातारे झाले आणि आपल्या तीन नातवंडांवर ---
दोन \VUgras{नातवांवर} \textsuperscript{(2)} आणि एका \VUgras{नातीवर}
\textsuperscript{(3)} --- फार प्रेम करू लागले. आज \VUgras{आजोबांचा वाढदिवस} आहे,
म्हणून मुले त्यांना शुभेच्छा द्यायला आली आहेत.

%%K t04-c2-02-1 p
२. --- लहानगा पीटर अजून फार लहान आहे; तो केवळ तीन वर्षांचा आहे. \VUgras{मांजर}
\textsuperscript{(1)} जवळ येताना तो पाहतो. तिच्या सुंदर रेशमी \VUgras{केसांवरून}
आणि तिच्या लांब \VUgras{शेपटीवरून} हात फिरवावा असे त्याला वाटते; त्या सुबक
\VUgras{पंजांखाली} टोकदार दुष्ट नखे लपली आहेत आणि ती कदाचित त्याला ओरबाडतील, याचा
तो विचार करत नाही.

%%K t04-c2-03-1 p
३. --- त्याचा हात धरणारी त्याची बहीण गॅब्रिएला पुष्कळ अधिक गंभीर आहे; आणि आपल्या
मोठ्या \VUgras{ओव्हरकोटासह} \textsuperscript{(4)} व कातडी \VUgras{पट्ट्यासह}
\textsuperscript{(5)} मार्सेलस जरा मोठा झाल्यासारखा दिसतो --- जरी त्याची आखूड
विजार त्याचे \VUgras{मोजे} \textsuperscript{(6)} दाखवत असली तरी.

%%K t04-c2-04-1 p
४. --- त्यांच्या वडिलांनी आपला जाड हिवाळी \VUgras{ओव्हरकोट}
\textsuperscript{(7)} घातला आहे --- \VUgras{फरच्या कॉलरचा} \textsuperscript{(8)}
--- आणि आपल्या \VUgras{खिशात} \textsuperscript{(9)} त्यांनी रुमाल ठेवला आहे.
त्यांच्या पत्नीच्या डोक्यावर \VUgras{पिसांनी} \textsuperscript{(10)} सजवलेली
लोकरी टोपी आहे, अंगावर आपले \VUgras{जाकीट} \textsuperscript{(11)}, गळ्यात आपला
\VUgras{फरचा बोआ} \textsuperscript{(12)} आणि हातांत आपला \VUgras{मफ}
\textsuperscript{(13)} आहे.

%%K t04-c2-05-1 p
५. --- फार थंडी आहे, कारण हिवाळ्याचे राज्य आहे. दिवसभर \VUgras{बर्फ}
\textsuperscript{(19)} पडत होता आणि घरांची छपरे अगदी पांढरी झाली आहेत.
\VUgras{रात्रीबरोबर} थंडी आणखी कडक होते आहे. आकाशात \VUgras{चंद्र}
\textsuperscript{(17)} आणि \VUgras{तारे} \textsuperscript{(18)} चमकत आहेत आणि
\VUgras{अंधार} भेदत आहेत.

%%K t04-tit-8 sub
\VUsaut{4.2mm}
\VUpk{10.2pt}{40}{आजार.}
\VUsaut{3.0mm}

%%K t04-c2-06-1 p
६. --- तो बिचारा \VUgras{आंधळा} \textsuperscript{(20)} फार दुर्दैवी आहे --- तो
\VUgras{औषधाच्या दुकानासमोरचे} \textsuperscript{(16)} चौक ओलांडतो आहे, आणि
त्याला त्याचा \VUgras{पुडल कुत्रा} \textsuperscript{(21)} वाट दाखवतो आहे.
जस्टिना, \VUgras{नोकराणी} \textsuperscript{(14)}, स्वयंपाकघरातून भरपूर साखर
घातलेल्या गरम \VUgras{काढ्याची} एक \VUgras{वाटी} \textsuperscript{(15)} आणते
आहे.

%%K t04-c2-07-1 p
७. --- \VUgras{आजी} \textsuperscript{(23)} टेबलावरचा आपला \VUgras{चष्मा}
\textsuperscript{(26)} शोधू पाहतात --- \VUgras{औषधाची चिठ्ठी}
\textsuperscript{(27)} वाचण्यासाठी, जी \VUgras{डॉक्टरांनी} \textsuperscript{(25)}
नुकतीच लिहिली आहे ---
आपल्या \VUgras{काठीवर} \textsuperscript{(24)} टेकून आरामखुर्चीतून उठतात.

%%K t04-c2-08-1 p
८. --- त्यांना वारंवार \VUgras{किरकोळ आजार}, \VUgras{चक्कर}, \VUgras{सर्दी} आणि
\VUgras{दातदुखी} होते; त्यांचे पाय अशक्त झाले आहेत आणि त्यांचे डोळेही आता फारसे
चांगले नाहीत. तरीही त्या आपल्या जुन्या \VUgras{डॉक्टरांची} थट्टा करायला मागेपुढे
पाहत नाहीत, जे त्यांना नेहमी \VUgras{औषध} \textsuperscript{(28)} लिहून द्यायला
तयार असतात. आज त्या फार आनंदी आहेत, कारण त्यांच्याभोवती त्यांचे तरुण कुटुंब आहे.

%%K t04-c2-09-1 p
९. --- नीट बंद केलेल्या खोलीत आराम वाटतो. संगमरवरी \VUgras{शेकोटीत}
\textsuperscript{(29)}

%%K t04-c2-09-1 p suite
\VUgras{सुंदर आग} \textsuperscript{(30)} जळते आहे, आणि नुकत्याच लावलेल्या
\VUgras{दिव्याच्या} \textsuperscript{(31)} मंद \VUgras{प्रकाशात} माणसाला सुख वाटते.

%%K t04-tit-9 sub
\VUsaut{4.2mm}
\VUpk{10.2pt}{40}{आजोबा.}
\VUsaut{3.0mm}

%%K t04-c2-10-1 p
१०. --- \VUgras{आजोबांनासुद्धा} \textsuperscript{(33)} विसर पडला आहे की ते आजारी
होते, की त्यांचा \VUgras{संधिवात} त्यांना त्यांच्या \VUgras{आरामखुर्चीला}
\textsuperscript{(34)} खिळवून ठेवतो, आणि की कालसुद्धा त्यांना \VUgras{ताप आणि
मूर्च्छा} आली होती. गेल्या हिवाळ्यात त्यांना \VUgras{न्यूमोनिया} झाला होता आणि
कोरड्या, त्रासदायक \VUgras{खोकल्याने} त्यांना फार त्रास दिला होता --- हेही त्यांना
आता आठवत नाही.

%%K t04-c2-10-2 p
पण त्यांनी \VUgras{उशीवर} \textsuperscript{(38)} टेकवलेला पायही त्यांना
दुखतो --- ती उशी एका लहान \VUgras{स्टुलावर} \textsuperscript{(39)} ठेवली आहे.

%%K t04-c2-11-1 p
११. --- ते आपल्या उबदार \VUgras{झग्यात} \textsuperscript{(35)} नीट लपेटलेले
असतात --- शिंपल्याच्या जाड \VUgras{बटणांच्या} \textsuperscript{(36)} झग्यात ---
आणि त्यांना वर्तमानपत्र वाचून दाखवायला जवळ ती लहानगी \VUgras{अनाथ मुलगी}
\textsuperscript{(40)} असते --- पांढरी टोपी, निळा \VUgras{झगा}
\textsuperscript{(42)} आणि \VUgras{शाल} \textsuperscript{(41)} घातलेली ---
तेव्हा ते तक्रार करत नाहीत.

%%K t04-c2-12-1 p
१२. --- दरवर्षी, जेव्हा \VUgras{दिनदर्शिका} \textsuperscript{(43)} पुन्हा एक
\VUgras{डिसेंबरची} \VUgras{तारीख} दाखवते, तेव्हा ते अधीरतेने आपल्या
\VUgras{नातवंडांची} वाट पाहतात, कारण त्यांना माहीत आहे की ती महत्त्वाची
\VUgras{तारीख} ती विसरणार नाहीत.

%%K t04-c2-13-1 p
१३. --- त्यांना भरून येऊन तो फार जुना काळ आठवतो, जेव्हा ते स्वतः जाऊन त्या
वैभवशाली राजेशाही \VUgras{सेनापतींना} शुभेच्छा देत असत, ज्यांची \VUgras{तसबीर}
\textsuperscript{(44)} भिंतीवर टांगलेली आहे आणि जे त्यांच्या मुलांचे
\VUgras{पणजोबा} आहेत.

\VUsaut{6.0mm}
\VUfilet{22mm}
